\subsection{terminal}\label{terminal}
Das \verb|terminal| Modul ist eine \verb|kybd_t|\-Implementierung
(Abschnitt \ref{keyboard}), die Tasteneingaben von einer seriellen Konsole
entgegennimmt. Es exportiert das Gerät \verb|terminal_dev| mit der
Betriebsart \verb|TERMINAL|. Intern hält es einen eigenen \verb|state_t|
(\verb|my_term|) mit acht Tasten, die mit \verb|'1'|\dots\verb|'8'|
beschriftet sind, und greift für den UART\-Zugriff auf die \verb|sio_t|\-Konfiguration
des \verb|serial|\-Moduls (Abschnitt \ref{serial}) zurück.

\subsubsection*{Funktionsweise}
{\sloppy\emergencystretch=3em
Beim Abtasten (\verb|terminal_scan|) fordert das Modul einmalig zur Eingabe
auf (\glqq Please enter key\grqq{}) und liest dann zeichenweise über
\verb|HAL_UART_Receive| bis zum Zeilenende. Eine erkannte Ziffer wird über
\verb|state_ch2idx()| auf einen Tastenindex abgebildet; liegt dieser im
gültigen Bereich, wird die Änderung in \verb|my_term| propagiert und das
\emph{dirty}\-Flag gesetzt. Die übrigen Vtable\-Methoden (\verb|terminal_state|,
\verb|terminal_diff|, \verb|terminal_add|, \verb|terminal_isdirty|, \dots)
greifen direkt auf die Funktionen des \verb|state|\-Moduls zu.
\par}

\subsubsection*{API}
{\sloppy\emergencystretch=3em
\begin{description}
\item[\texttt{terminal\_waitForNumber(key)}] Blockierend: liest eine Zahl von
  der seriellen Konsole; \verb|key| zeigt danach auf die gelesene
  Zeichenkette. Die Eingabe \verb|R|/\verb|r| signalisiert einen Reset.
  Rückgabe: \verb|int8_t| -- der Wert, bzw.\ $-1$ bei Reset/ungültiger
  Eingabe.
\end{description}
\par}

\noindent Daneben existiert ein \verb|UART_IdleCallback()| für den
DMA\-basierten Empfang (Idle\-Line\-Erkennung), der empfangene Zeichen in das
Label \verb|clabel| des \verb|state_t| übernimmt.
